\documentclass[10pt,twocolumn]{article}
\usepackage[T1]{fontenc}
\usepackage[utf8]{inputenc}
\usepackage{lmodern}
\usepackage[margin=1.8cm,top=2.4cm,bottom=2cm,columnsep=0.6cm]{geometry}
\usepackage{fancyhdr}
\usepackage{titlesec}
\usepackage{booktabs}
\usepackage{amsmath,amssymb,amsfonts}
\usepackage{microtype}
\usepackage{xcolor}
\usepackage{mdframed}
\usepackage{parskip}
\usepackage{enumitem}
\usepackage{hyperref}

\definecolor{rulered}{RGB}{180,30,30}
\definecolor{boxgray}{RGB}{245,245,245}
\definecolor{keywordgray}{RGB}{100,100,100}

\pagestyle{fancy}
\fancyhf{}
\fancyhead[L]{\textsc{\small Claude Mag}}
\fancyhead[C]{\small\textit{Machine Learning Research Digest}}
\fancyhead[R]{\small 2026-07-06}
\fancyfoot[C]{\small\thepage}
\renewcommand{\headrulewidth}{0.8pt}

\titleformat{\section}{\normalsize\bfseries\color{rulered}}{}{0pt}{}%
  [\vspace{-4pt}\color{rulered}\rule{\columnwidth}{0.4pt}]
\titlespacing{\section}{0pt}{8pt}{4pt}

\hypersetup{colorlinks=true,urlcolor=rulered,linkcolor=black}

\begin{document}
\setlength{\parindent}{0pt}
\setlength{\parskip}{4pt}

{\small\color{keywordgray}\textbf{Keywords:}
  multi-turn conversation \textbullet{}
  LLM evaluation \textbullet{}
  performance degradation \textbullet{}
  conversational AI}\par\vspace{4pt}

{\LARGE\bfseries\raggedright
  LLMs Get Lost In Multi-Turn Conversation}\par\vspace{4pt}

{\large\itshape\raggedright
  All Frontier Models Suffer a 39\% Average Performance Drop
  Moving from Single-Turn to Multi-Turn Settings}\par\vspace{6pt}

{\small\textbf{By} Philippe Laban, Hiroaki Hayashi, Yingbo Zhou \&
  Jennifer Neville (Salesforce AI Research)}\par\vspace{2pt}

{\small\color{keywordgray}arXiv:2505.06120 \textbullet{}
  ICLR 2026 Outstanding Paper Award}
\par\vspace{2pt}\noindent{\color{rulered}\rule{\columnwidth}{0.6pt}}\vspace{6pt}

Every frontier language model tested---open- and closed-weight---performs
significantly worse in multi-turn conversations than in equivalent single-turn
tasks. Across more than 200,000 simulated conversations spanning six generation
tasks, the average performance drop is \textbf{39\%}. More troublingly,
once a model takes a wrong conversational turn, it fails to recover:
it gets lost and stays lost.

\begin{mdframed}[backgroundcolor=boxgray,linecolor=rulered,linewidth=1pt,
    innerleftmargin=6pt,innerrightmargin=6pt,
    innertopmargin=6pt,innerbottommargin=6pt]
\textbf{\small AT A GLANCE}\par\vspace{4pt}
\begin{description}[leftmargin=0pt,labelsep=4pt,itemsep=2pt,parsep=0pt,
    font=\small\bfseries]
  \item[Problem:] LLMs are evaluated almost exclusively in single-turn,
    fully-specified instruction settings, masking a critical failure mode in
    multi-turn dialogue.
  \item[Why it matters:] Real-world use involves extended conversations where
    user intent is underspecified and refined iteratively; models that degrade
    are unsafe to deploy in agentic contexts.
  \item[Method:] 200,000+ simulated conversations on six generation tasks,
    decomposing the performance gap into aptitude loss and misalignment loss.
  \item[Result:] 39\% average drop across all frontier models; recovery after a
    wrong turn is essentially zero across all settings.
  \item[Evidence:] GPT-4o, Claude 3.5 Sonnet, Gemini 1.5 Pro, Llama 3.1
    (70B/405B), Mistral Large all exhibit the same pattern.
\end{description}
\end{mdframed}

\section{The Big Picture}

LLMs are conversational interfaces by design, yet nearly every benchmark
measures performance on single isolated turns with fully specified prompts.
This gap matters: analysis of real user logs shows that underspecification is
the norm, not the exception. Users iterate, refine, and clarify. The question
of how models perform across that iterative process has been largely ignored.

This paper closes that gap with a large-scale simulation study comparing
single-turn and multi-turn performance on the same underlying tasks. The
results are stark and consistent: every model tested degrades substantially
in the multi-turn setting, and none recovers once a conversational wrong
turn has been taken.

\section{Why Existing Approaches Fall Short}

Standard benchmarks like MMLU, HumanEval, and MT-Bench evaluate models
on isolated, well-formed prompts. Even ``multi-turn'' benchmarks typically
measure whether a model can carry context forward, not whether it can
integrate iterative user refinements without losing the original intent.

The underspecification problem is orthogonal to context length: even short
conversations (3--5 turns) show significant degradation. This rules out
context-length limitations as the primary cause and points instead to a
failure in \emph{intent tracking} as conversation state evolves.

\section{Core Mechanism}

The authors decompose the multi-turn performance drop into two components:

\textbf{Aptitude loss:} A minor degradation in the model's core task ability
when presented with conversation history rather than a direct prompt---small
but measurable ($\approx 5$--8 percentage points).

\textbf{Misalignment loss:} The dominant effect ($\approx 31$--35 percentage
points on average). As the conversation progresses, the model increasingly
anchors to its own prior outputs rather than the user's current intent.
A single wrong inference early in the conversation corrupts all downstream
turns---the ``getting lost'' phenomenon.

\textbf{Non-recovery:} The recovery rate $R_n$---the probability that
performance at turn $n+k$ exceeds performance at turn $n$ after a detected
misalignment---is approximately zero across all models, tasks, and
conversation lengths tested.

\section{Technical Formulation}

Let $T = (t_1, t_2, \ldots, t_n)$ be a conversation and $y$ the target
output. Model performance at turn $n$ is:
\[
\text{Perf}_n = P\!\left(f(T_{1:n}) = y\right).
\]
The paper observes $\text{Perf}_n < \text{Perf}_1$ for all $n > 1$, all
models, all tasks.

The \emph{recovery rate} is defined as:
\[
R_n = \frac{1}{K}\sum_{k=1}^{K}
  \mathbf{1}\!\left[\text{Perf}_{n+k} > \text{Perf}_n\right].
\]
Empirically, $R_n \approx 0$ across all conditions.

The performance gap is decomposed as:
\[
\Delta = \underbrace{(\text{Perf}_1 - \text{Perf}_1^{\text{hist}})}_{\text{aptitude loss}}
       + \underbrace{(\text{Perf}_1^{\text{hist}} - \text{Perf}_n)}_{\text{misalignment loss}}
\]
where $\text{Perf}_1^{\text{hist}}$ is single-turn performance when history
is present but the core prompt is unchanged.

\section{Learning or Inference Procedure}

The evaluation protocol creates multi-turn conversations by simulation:
\begin{enumerate}[itemsep=2pt,parsep=0pt]
  \item Sample a task instance $(x, y)$ from each benchmark.
  \item Construct an underspecified initial prompt $t_1$ that omits key
    constraints or details.
  \item Simulate a user who provides clarifications at each turn $t_k$,
    moving toward the fully-specified target.
  \item Score the model's response at each turn against $y$.
  \item Aggregate across 200,000+ conversations.
\end{enumerate}

\section{What the Guarantee Says}

The paper does not provide formal guarantees, but its empirical finding
is stated as a universal claim: \emph{no evaluated model} achieves
multi-turn performance within 20\% of its single-turn performance,
and \emph{no evaluated model} has a recovery rate above 5\% after
a wrong turn.

\section{Experimental Findings}

Evaluated on six tasks---summarization, email drafting, code generation,
QA, creative writing, data extraction---the paper reports:
\begin{itemize}[itemsep=2pt,parsep=0pt]
  \item 39\% average performance drop across all models and tasks.
  \item Code generation and QA: largest drops ($\approx$45\%).
  \item Creative writing: smallest drop ($\approx$28\%).
  \item Closed-weight models (GPT-4o, Claude, Gemini) do not outperform
    open-weight models (Llama, Mistral) on multi-turn robustness.
  \item Performance degradation begins at turn 2 and worsens monotonically
    through turn 5, then plateaus.
\end{itemize}

\section{Ablations and Interpretation}

\textbf{Turn count:} Degradation is sharpest between turns 1 and 2, then
increases more slowly. The model's initial interpretation of the underspecified
prompt is highly persistent.

\textbf{Clarification type:} Constraints added at turn 3+ are less reliably
incorporated than those present at turn 1. Retroactive clarifications are
especially poorly handled.

\textbf{Model scale:} Larger models within the same family show slightly
smaller aptitude loss but equivalent misalignment loss --- scale alone does
not solve multi-turn robustness.

\textbf{Implications:} The result motivates dedicated multi-turn training
objectives, retrieval-augmented intent tracking, and explicit user-intent
summarization as architectural additions.

\section*{Reference}

Philippe Laban, Hiroaki Hayashi, Yingbo Zhou, Jennifer Neville.
\textbf{LLMs Get Lost In Multi-Turn Conversation}.
arXiv:2505.06120, May 2025.
\url{https://arxiv.org/abs/2505.06120}

\end{document}
